% GENERATED FILE. Edit canonical lessons, then run npm run generate:books.
% canonical-source-hash: fnv1a64:79eebd29a1d3b63f
% canonical-lessons: KA-C498-nambu, KA-C498-mannisu, KA-C498-samsaya, KA-C498-karune, KA-C498-sahanubhuti

\chapter{Trust and Sympathy}
\label{ch:ka-trust-and-sympathy}
\hlchaptermodality{498}

\begin{chapteropening}
\textbf{By the end of this chapter:} \emph{I can say to believe, to forgive, suspicion, compassion and sympathy.}

Complete the last lesson of chapter 498: I can say to believe, to forgive, suspicion, compassion and sympathy.
\end{chapteropening}

\section[\texorpdfstring{\kn{ನಂಬು} (nambu)}{ನಂಬು (nambu)}]{\kn{ನಂಬು} (nambu) — to believe; to trust}

\label{lesson:KA-C498-nambu}



\begin{quote}
Before the new one: say the Kannada for modesty, then the Kannada for self-respect.
\end{quote}

\subsection*{You'll want to know: \kn{ನಂಬು}}
\textbf{\kn{ನಂಬು}} — \emph{nambu} — \textquotedblleft{}to believe; to trust\textquotedblright{}.

\begin{quote}
\kn{ಅವನ} \kn{ಮಾತನ್ನು} \kn{ಯಾರೂ} \kn{ನಂಬಲಿಲ್ಲ}.
\end{quote}

\begin{quote}
\emph{avana mātannu yārū nambalilla.}
\end{quote}

\begin{quote}
Nobody believed what he said.
\end{quote}

Past (I …): \kn{ನಂಬಿದೆ}, \emph{nambide}.

What or whom you believe takes -\kn{ಅನ್ನು}. The noun is \kn{ನಂಬಿಕೆ}, trust.

\begin{cousinweb}
\textbf{\kn{ನಂಬು}} (\emph{nambu}) is native Dravidian: DEDR 3600 lists Kannada \emph{nambu}, \textquotedblleft{}to trust, believe\textquotedblright{}, and \emph{nambike}, \textquotedblleft{}trust\textquotedblright{}, beside Tamil \emph{nampu}.
\end{cousinweb}

\begin{grammarlens}[title={What you've built}]
One more word for describing what you lived and felt.
\end{grammarlens}

\subsection*{Your turn}
\begin{itemize}
\raggedright
  \item \emph{Say it:} \emph{nambu}
  \item \emph{Say it:} \emph{nambu}, once more
  \item \emph{Say it:} \emph{avana mātannu yārū nambalilla.} — Nobody believed what he said.
  \item \emph{Recall:} say the Kannada for modesty, then the Kannada for self-respect, then say \emph{nambu} again
\end{itemize}

\begin{tcolorbox}[breakable,colback=teal!4,colframe=teal!35!black,title={Before you move on}]
What does \kn{ನಂಬು} mean? (\textquotedblleft{}To believe; to trust\textquotedblright{}.) Say it once more.
\end{tcolorbox}
\section[\texorpdfstring{\kn{ಮನ್ನಿಸು} (mannisu)}{ಮನ್ನಿಸು (mannisu)}]{\kn{ಮನ್ನಿಸು} (mannisu) — to forgive; to excuse}

\label{lesson:KA-C498-mannisu}



\begin{quote}
Before the new one: say the Kannada for self-respect, then the Kannada for to believe.
\end{quote}

\subsection*{You'll want to know: \kn{ಮನ್ನಿಸು}}
\textbf{\kn{ಮನ್ನಿಸು}} — \emph{mannisu} — \textquotedblleft{}to forgive; to excuse\textquotedblright{}.

\begin{quote}
\kn{ತಡವಾಗಿ} \kn{ಬಂದಿದ್ದಕ್ಕೆ} \kn{ಮನ್ನಿಸಿ} \kn{ಅಂತ} \kn{ಅವನು} \kn{ಕೇಳಿಕೊಂಡನು}.
\end{quote}

\begin{quote}
\emph{taḍavāgi bandiddakke mannisi anta avanu kēḷikoṇḍanu.}
\end{quote}

\begin{quote}
He asked to be excused for coming late.
\end{quote}

Past (I …): \kn{ಮನ್ನಿಸಿದೆ}, \emph{manniside}.

Polite: \kn{ಮನ್ನಿಸಿ}, excuse me, close to \kn{ಕ್ಷಮಿಸಿ}.

\begin{cousinweb}
\textbf{\kn{ಮನ್ನಿಸು}} (\emph{mannisu}) is usually traced to Sanskrit \emph{mānana}, \textquotedblleft{}honouring, showing respect\textquotedblright{} (Monier-Williams), with the verb ending -\emph{isu}; to pardon someone is to treat them with regard. The path is not settled.
\end{cousinweb}

\begin{grammarlens}[title={What you've built}]
One more word for describing what you lived and felt.
\end{grammarlens}

\subsection*{Your turn}
\begin{itemize}
\raggedright
  \item \emph{Say it:} \emph{mannisu}
  \item \emph{Say it:} \emph{mannisu}, once more
  \item \emph{Say it:} \emph{taḍavāgi bandiddakke mannisi anta avanu kēḷikoṇḍanu.} — He asked to be excused for coming late.
  \item \emph{Recall:} say the Kannada for self-respect, then the Kannada for to believe, then say \emph{mannisu} again
\end{itemize}

\begin{tcolorbox}[breakable,colback=teal!4,colframe=teal!35!black,title={Before you move on}]
What does \kn{ಮನ್ನಿಸು} mean? (\textquotedblleft{}To forgive; to excuse\textquotedblright{}.) Say it once more.
\end{tcolorbox}
\section[\texorpdfstring{\kn{ಸಂಶಯ} (saṃśaya)}{ಸಂಶಯ (saṃśaya)}]{\kn{ಸಂಶಯ} (saṃśaya) — suspicion; doubt}

\label{lesson:KA-C498-samsaya}



\begin{quote}
Before the new one: say the Kannada for to believe, then the Kannada for to forgive.
\end{quote}

\subsection*{You'll want to know: \kn{ಸಂಶಯ}}
\textbf{\kn{ಸಂಶಯ}} — \emph{saṃśaya} — \textquotedblleft{}suspicion; doubt\textquotedblright{}.

\begin{quote}
\kn{ಆ} \kn{ಅಪರಿಚಿತನ} \kn{ಮೇಲೆ} \kn{ಪೊಲೀಸರಿಗೆ} \kn{ಸಂಶಯ} \kn{ಬಂದಿತು}.
\end{quote}

\begin{quote}
\emph{ā aparicitana mēle polīsarige saṃśaya banditu.}
\end{quote}

\begin{quote}
The police grew suspicious of the stranger.
\end{quote}

On whom: … \kn{ಮೇಲೆ} \kn{ಸಂಶಯ}. Close to \kn{ಅನುಮಾನ}.

\begin{cousinweb}
\textbf{\kn{ಸಂಶಯ}} (\emph{saṃśaya}) is borrowed from Sanskrit \emph{saṃśaya}, \textquotedblleft{}doubt, uncertainty\textquotedblright{} (Monier-Williams).
\end{cousinweb}

\begin{grammarlens}[title={What you've built}]
One more word for describing what you lived and felt.
\end{grammarlens}

\subsection*{Your turn}
\begin{itemize}
\raggedright
  \item \emph{Say it:} \emph{saṃśaya}
  \item \emph{Say it:} \emph{saṃśaya}, once more
  \item \emph{Say it:} \emph{ā aparicitana mēle polīsarige saṃśaya banditu.} — The police grew suspicious of the stranger.
  \item \emph{Recall:} say the Kannada for to believe, then the Kannada for to forgive, then say \emph{saṃśaya} again
\end{itemize}

\begin{tcolorbox}[breakable,colback=teal!4,colframe=teal!35!black,title={Before you move on}]
What does \kn{ಸಂಶಯ} mean? (\textquotedblleft{}Suspicion; doubt\textquotedblright{}.) Say it once more.
\end{tcolorbox}
\section[\texorpdfstring{\kn{ಕರುಣೆ} (karuṇe)}{ಕರುಣೆ (karuṇe)}]{\kn{ಕರುಣೆ} (karuṇe) — compassion; mercy}

\label{lesson:KA-C498-karune}



\begin{quote}
Before the new one: say the Kannada for to forgive, then the Kannada for suspicion.
\end{quote}

\subsection*{You'll want to know: \kn{ಕರುಣೆ}}
\textbf{\kn{ಕರುಣೆ}} — \emph{karuṇe} — \textquotedblleft{}compassion; mercy\textquotedblright{}.

\begin{quote}
\kn{ಬಡವರ} \kn{ಮೇಲೆ} \kn{ಅವಳಿಗೆ} \kn{ತುಂಬಾ} \kn{ಕರುಣೆ} \kn{ಇದೆ}.
\end{quote}

\begin{quote}
\emph{baḍavara mēle avaḷige tumbā karuṇe ide.}
\end{quote}

\begin{quote}
She has great compassion for the poor.
\end{quote}

Close to \kn{ದಯೆ}, kindness; \kn{ಕರುಣೆ} is pity that moves you.

\begin{cousinweb}
\textbf{\kn{ಕರುಣೆ}} (\emph{karuṇe}) is borrowed from Sanskrit \emph{karuṇā}, \textquotedblleft{}pity\textquotedblright{} (CDIAL 2811); the final -ā became -\kn{ೆ}, as in many Kannada loans.
\end{cousinweb}

\begin{grammarlens}[title={What you've built}]
One more word for describing what you lived and felt.
\end{grammarlens}

\subsection*{Your turn}
\begin{itemize}
\raggedright
  \item \emph{Say it:} \emph{karuṇe}
  \item \emph{Say it:} \emph{karuṇe}, once more
  \item \emph{Say it:} \emph{baḍavara mēle avaḷige tumbā karuṇe ide.} — She has great compassion for the poor.
  \item \emph{Recall:} say the Kannada for to forgive, then the Kannada for suspicion, then say \emph{karuṇe} again
\end{itemize}

\begin{tcolorbox}[breakable,colback=teal!4,colframe=teal!35!black,title={Before you move on}]
What does \kn{ಕರುಣೆ} mean? (\textquotedblleft{}Compassion; mercy\textquotedblright{}.) Say it once more.
\end{tcolorbox}
\section[\texorpdfstring{\kn{ಸಹಾನುಭೂತಿ} (sahānubhūti)}{ಸಹಾನುಭೂತಿ (sahānubhūti)}]{\kn{ಸಹಾನುಭೂತಿ} (sahānubhūti) — sympathy}

\label{lesson:KA-C498-sahanubhuti}



\begin{quote}
Before the new one: say the Kannada for suspicion, then the Kannada for compassion.
\end{quote}

\subsection*{You'll want to know: \kn{ಸಹಾನುಭೂತಿ}}
\textbf{\kn{ಸಹಾನುಭೂತಿ}} — \emph{sahānubhūti} — \textquotedblleft{}sympathy\textquotedblright{}.

\begin{quote}
\kn{ನೆರೆಯವರು} \kn{ನಮ್ಮ} \kn{ನಷ್ಟಕ್ಕೆ} \kn{ಸಹಾನುಭೂತಿ} \kn{ತೋರಿಸಿದರು}.
\end{quote}

\begin{quote}
\emph{nereyavaru namma naṣṭakke sahānubhūti tōrisidaru.}
\end{quote}

\begin{quote}
The neighbours showed sympathy for our loss.
\end{quote}

\kn{ಸಹ}, with + \kn{ಅನುಭೂತಿ}, feeling: feeling along with someone.

\begin{cousinweb}
\textbf{\kn{ಸಹಾನುಭೂತಿ}} (\emph{sahānubhūti}) is a learned compound built from Sanskrit \emph{saha}, \textquotedblleft{}with\textquotedblright{}, and \emph{anubhūti}, \textquotedblleft{}perception, experience\textquotedblright{} (Monier-Williams).
\end{cousinweb}

\begin{grammarlens}[title={What you've built}]
One more word for describing what you lived and felt.
\end{grammarlens}

\subsection*{Your turn}
\begin{itemize}
\raggedright
  \item \emph{Say it:} \emph{sahānubhūti}
  \item \emph{Say it:} \emph{sahānubhūti}, once more
  \item \emph{Say it:} \emph{nereyavaru namma naṣṭakke sahānubhūti tōrisidaru.} — The neighbours showed sympathy for our loss.
  \item \emph{Recall:} say the Kannada for suspicion, then the Kannada for compassion, then say \emph{sahānubhūti} again
\end{itemize}

\begin{tcolorbox}[breakable,colback=teal!4,colframe=teal!35!black,title={Before you move on}]
What does \kn{ಸಹಾನುಭೂತಿ} mean? (\textquotedblleft{}Sympathy\textquotedblright{}.) Say it once more.
\end{tcolorbox}
